% Meta-Diffusion: what has been run, and what each run established
%
% Compile:  pdflatex results_summary.tex
% Requires: amsmath amssymb booktabs enumitem xcolor hyperref geometry
%
% Every figure quoted here was read out of checkpoints/*_log.jsonl or artifacts/ on
% 11 September 2026, not from recollection. Where a number is a single diagnostic point
% rather than an average, it says so.

\documentclass[11pt]{article}

\usepackage[a4paper,margin=1in]{geometry}
\usepackage{amsmath,amssymb}
\usepackage{booktabs}
\usepackage{enumitem}
\usepackage{xcolor}
\usepackage{hyperref}

\definecolor{shared}{HTML}{1D4E4A}
\definecolor{overturn}{HTML}{8C3A32}

\hypersetup{colorlinks=true, linkcolor=shared, urlcolor=shared}

\newcommand{\code}[1]{\texttt{\small #1}}
\newcommand{\res}{\par\smallskip\textit{Result.}\ }
\newcommand{\shows}{\par\smallskip\textit{Shows.}\ }
\newcommand{\dtask}{\Delta_{\mathrm{task}}}
\newcommand{\KT}{K_{T}}

\setlength{\emergencystretch}{3em}
\setlist[enumerate]{leftmargin=1.8em, itemsep=8pt, topsep=6pt}

\title{\vspace{-1.5em}Meta-Diffusion: experiments run so far, and what each established}
\date{11 September 2026}
\author{}

\begin{document}
\maketitle
\vspace{-2em}

\noindent
The method trains one diffusion network across many tasks, freezes it, and adapts to a new
task by fitting a $k$-dimensional coordinate $z$ that weights $k$ learned
score-perturbation directions:
\[
  \hat\varepsilon(x_t,t,z) \;=\; \hat\varepsilon_0(x_t,t) \;+\; \textstyle\sum_{\ell} z_\ell\,R_\ell(x_t,t).
\]
A transport map predicts the target coordinate from source data alone. Eleven experiments
have been run. Two domains were used: two-dimensional Gaussian mixtures, where the true
score is analytic, and CIFAR-100, where it is not.

\section{What has been run}

\begin{enumerate}

\item \textbf{Analytic validation on two-dimensional Gaussian mixtures (E1).}
Four-component mixtures, where the true score has a closed form, so model output is
compared against truth rather than against sample appearance.
\res Source-informed transport beats every baseline at $\KT = 1$, and the advantage
shrinks as target evidence grows, with paired confidence intervals over 32 tasks and 3
seeds.
\shows The intended mechanism is real and the implementation can express it. This remains
the strongest positive evidence in the programme.

\item \textbf{CIFAR-100 with sibling fine classes as source and target (E2).}
A deliberately hard, low-relatedness stress test.
\res The coordinate collapses: relative spread across episodes stays near $0.02$ against a
$0.05$ threshold. Dissecting a checkpoint showed the pooled encoder features barely
separate one flower species from another.
\shows The encoder cannot extract task identity from the denoising gradient alone when
tasks differ only by a sibling class. The protocol, not the model, needed changing.

\item \textbf{CIFAR-100 with one fixed transformation, clean to blurred (E3).}
The corrected protocol: same semantic class, clean source images, transformed target
images, superclasses split 12/4/4 into train, validation and test.
\res Collapse is solved --- coordinate spread rises to $0.36$ --- but the benefit
disappears. Basis usage $r_{\mathrm{basis}}$ falls from $0.254$ to $0.011$, and the gain
over $z = 0$ peaks at $0.084$ near step 2000 and reaches $0.0002$ by step 50{,}000.
\shows As the 27.8M-parameter backbone matures it models clean and blurred images
unconditionally, leaving the coordinate nothing to do.

\item \textbf{CIFAR-100 with three transformations: blur, noise, contrast (E4).}
The hypothesis was that a backbone unable to know \emph{which} corruption applies would be
forced to use the coordinate.
\res It was not. The same arc repeats: gain peaks at $0.078$ and settles at $0.0015$;
$r_{\mathrm{basis}}$ ends at $0.037$. The encoder meanwhile keeps improving, reaching a
coordinate spread of $1.19$.
\shows Encoder quality and coordinate usefulness are decoupled. A coordinate can become
more informative while contributing less.

\item \textbf{Coordinate dimension sweep, $k \in \{16, 32, 64\}$ (E5).}
Trained separately at matched steps, measured on validation.
\res A larger $k$ raises the early peak --- $0.084$, $0.105$, $0.129$ --- and then decays
slightly further, to $0.0007$, $0.0005$ and $0.0004$ at step 20{,}000, with all three
curves still falling.
\shows $k$ is not the bottleneck. A larger basis is absorbed faster, which is the
signature of a backbone taking over the difference rather than of a coordinate short of
capacity.

\item \textbf{Backbone capacity sweep, 128 / 64 / 32 channels at fixed $k = 32$ (E6).}
\res Shrinking the backbone changes the picture completely at 20{,}000 steps: basis usage
rises from $0.024$ to $0.67$ to $1.37$, and the gain over $z=0$ from $0.0005$ to $0.163$
to $0.386$ --- a 725-fold rise.
\shows \textcolor{overturn}{This was read as the mechanism switching on, and that reading
was wrong.} What grew was the size of a constant offset; see E9.

\item \textbf{Convergence of the two small backbones to 60{,}000 steps (E7).}
\res Every curve is still falling, with the half-life of an exponential fit scaling with
capacity: 1{,}965 steps at 128 channels, 10{,}160 at 64 and 25{,}380 at 32. Meanwhile the
denoising loss under the correct coordinate stays flat throughout.
\shows Capacity buys time, not permanence, and this is not undertraining --- the model has
converged while the coordinate keeps losing relevance.

\item \textbf{Deployment sweeps over $\KT \in \{1,2,5,10,20\}$ and source evidence (E8).}
Six adaptation strategies on 40 held-out episodes, paired per episode.
\res Any non-zero coordinate beats $z = 0$ by $0.0015$ (95\,\% CI $\pm 0.0002$) ---
statistically real, but 1.7\,\% of the loss. \emph{No coordinate beats any other}:
transport, refinement and the oracle upper bound all land on $0.0887$, and $\KT$ changes
nothing.
\shows Whatever the coordinate contributes does not depend on which coordinate it is.

\item \textbf{Mean-coordinate substitution (E9).} \emph{The decisive experiment.}
Replace each episode's own coordinate with a single mean coordinate shared by every
episode, which removes all task-specific content while keeping the common offset.
\res On the checkpoint where the apparent gain looked healthiest: own coordinate
$0.0273$, one shared mean $0.0273$, $z = 0$ giving $0.1456$.
\shows \textbf{Zero per cent of the coordinate's benefit is task-specific.} Its whole
contribution is a constant the model needs; $z = 0$ is catastrophic only because it
deletes that constant. This also re-reads E6: a large backbone absorbs the constant into
its weights, a small one leaves it in the basis, and neither adapts.

\item \textbf{Loss sensitivity to coordinate error (E10).}
\res The denoising loss around the target coordinate is very flat, and every coordinate
the pipeline produces falls inside that basin.
\shows Substantial coordinate error costs almost nothing, so coordinate-space accuracy
cannot be inferred from the loss, nor the loss from it.

\item \textbf{Coordinate accuracy against an abundant-target reference (E11).}
Encode a reference from 230 target images, then measure how far the source coordinate sits
from it before and after transport. 40 validation episodes.
\res Transport closes \textbf{93\,\% of the distance}, from $3.106 \pm 0.073$ to $0.227
\pm 0.019$. At $\KT = 1$ the transported coordinate is closer to the truth ($0.227$) than
the coordinate encoded from the single available target image ($0.397$); abundant evidence
overtakes it between five and twenty samples ($0.180$, $0.112$).
\shows The encoder reads tasks and transport predicts them, both demonstrably. The failure
is downstream of both: the basis cannot turn a correct coordinate into a denoising gain.

\end{enumerate}

\section{Added in the current round, from re-reading the logs}

\begin{enumerate}[resume]

\item \textbf{The shuffled control is not zero on the three-transformation arm.}
\res Averaged over the last ten diagnostic points of that run, substituting \emph{another
episode's} coordinate costs $+0.0029$, while deleting the coordinate entirely costs
$+0.0016$. On every other run this quantity is $+0.00003$.
\shows A wrong coordinate is worse than no coordinate, but only where episodes differ in
which transformation applies. The working hypothesis this raises is that the coordinate
carries transformation identity, not class identity --- which would explain why class-only
task families produce nothing.

\item \textbf{No training run on disk records the decisive control.}
\res The mean-coordinate substitution was added to the code after every CIFAR run had
already finished, so none of the eighteen logs contains it. E9 is a single post-hoc
measurement on one checkpoint, not a curve against training step.
\shows The quantity the next stage turns on has never been watched during training. It
cannot be recovered from existing data, and requires new runs whatever else is decided.

\item \textbf{Pairing is worth a factor of sixty-five.}
\res Over the converged stretch of one run, the standard deviation of the raw denoising
loss between diagnostic points is $0.0241$, while the standard deviation of the paired
difference is $0.00037$. The effect being measured is $0.0015$.
\shows The effect is sixteen times smaller than the between-point noise it sits in. A raw
loss curve cannot show it, and any comparison that is not paired returns an arbitrary sign.

\end{enumerate}

\section{Where this leaves the project}

Stated so it can be falsified: \textbf{the method needs a task family whose members lie
farther apart than the shared backbone can absorb.} Stage 1 qualifies, because one
unconditional model cannot cover four mixtures with modes in genuinely different places.
CIFAR-100 class-conditional distributions do not --- denoising a $32 \times 32$ natural
image scarcely depends on knowing its class, so ignoring $z$ is the correct solution to the
problem as posed.

Three of the four stages of the pipeline are therefore proven to work: the encoder reads a
task (E11), transport predicts one (E11), and the mechanism delivers when tasks are far
enough apart (E1). The fourth --- the basis turning a correct coordinate into a
task-specific change in the reverse dynamics --- has never been observed on images (E9).

Two rounds of experiments were reported as successes on the strength of the wrong number
and had to be withdrawn (E6, E7). The control that caught it costs one extra forward pass,
and is now computed on every run. The order of reading matters: the fraction of the
benefit that is task-specific, then the gain over $z = 0$, then basis usage.

\paragraph{The open question.}
\[
\boxed{\text{Does the task-specific part of } z \text{ produce a task-specific change in the reverse dynamics?}}
\]
On the denoising loss the answer so far is no. At the level of generated samples it has
never been measured, and that measurement is the first item of the next stage.

\end{document}
